\Section{Множества и кривые в \(\Co\)}

\begin{definition}
\textit{\(\epsilon\)-окрестностью точки \(z_0\)} будем называть множество \(B_\epsilon(z_0) = \{z
\in \Co \colon |z - z_0| < \epsilon\}\).
\end{definition}

\begin{definition}
\textit{Бесконечно удалённую точку на комплексной плоскости} будем обозначать знаком \(\infty\).
\(\overline{\Co} = \Co \cup \{\infty\}\) - расширенная комплексная плоскость.
\end{definition}

В случае \(\epsilon\)-окрестности бесконечно удалённой точки имеем дело с множеством \(B_\epsilon
(\infty) = \{z \in \Co \colon |z| > \epsilon\} \cup \{\infty\}\).

\begin{definition}
\textit{Проколотой \(\epsilon\)-окрестностью точки \(z_0 \in \overline{\Co}\)} будем называть
множество \(\mathring{B}_\epsilon(z_0) = B_\epsilon(z_0) \setminus \{z_0\}\).
\end{definition}

\begin{definition}
\textit{Множество \(D \subset \oCo\)} называется \textit{открытым}, если \(\forall z_0 \in D \
\exists \epsilon > 0 \colon B_\epsilon(z_0) \subset D\).
\end{definition}

\begin{definition}
\textit{Множество \(D \subset \oCo\)} называется \textit{ограниченным}, если \(\exists R > 0 \colon
D \subset B_R(0)\).
\end{definition}

\begin{definition}
\textit{Точка \(z_0 \in \oCo\)} называется \textit{предельной} для множества \(E\), если \(\forall
\epsilon > 0\) \(E \cap B_\epsilon(z_0)\) содержат бесконечное число точек. Стоит отметить, что
граничные точки могут как принадлежать, так и не принадлежать множеству.
\end{definition}

\begin{definition}
\textit{Точка \(z_0 \in \oCo\)} называется \textit{граничной} для множества \(D\), если \(\forall
\epsilon > 0\) \(B_\epsilon(z_0)\) содержат точки как из \(D\), так и вне \(D\).
\end{definition}

\begin{definition}
\textit{Границей множества \(D\)} будем называть множество \(\partial D\), состоящее из всех
граничных точек множества \(D\).
\end{definition}

\begin{definition}
\textit{Множество \(D \subset \oCo\)} называется \textit{замкнутым}, если оно содержит все свои
граничные точки.
\end{definition}

Перейдём к рассмотрению кривых на комплексной плоскости.

\begin{definition}
\textit{Кривой \(\gamma\)} будем называть непрерывное отображение некоторого отрезка \([a, b]\) в
\(\Co\). Выражение \(z = z(t) = x(t) + iy(t)\), где \(x, y\)~--- непрерывные на \([a,b]\) функции,
будем называть \textit{параметрическим отображением \(\gamma\)}. Кривую, с ориентацией обратной
ориентации \(\gamma\) будем обозначать \(\gamma^{-1}\), ей характерна параметризация \(z = z(-t)\).
Составную кривую \(\gamma\), состоящую из кривых \(\gamma_1, \dots, \gamma_n\) будем обозначать
\(\gamma = \gamma_1 \gamma_2 \dots \gamma_n\).
\end{definition}

\begin{definition}
\textit{Кривую \(\gamma\)} будем называть \textit{замкнутой} или \textit{контуром}, если \(z(a) =
z(b)\).
\end{definition}

\begin{definition}
\textit{Незамкнутую кривую \(\gamma\)} будем называть \textit{простой (жордановой)}, если отображение
\(z(t)\) инъективно.
\end{definition}

\begin{definition}
\textit{Замкнутую кривую \(\gamma\)} будем называть \textit{простой}, если сужение отображения
\(z(t)\) на \([a, b)\) инъективно.
\end{definition}

\begin{definition}
\textit{Кривую \(\gamma\)} будем называть \textit{непрерывно-дифференцируемой}, если функции \(x, y\)
в отображении \(z\) непрерывно-дифференцируемы. Если при этом они одновременно не обращаются в нуль
\(\forall t \in [a,b]\), то \textit{кривую} будем называть \textit{гладкой}.
\end{definition}

\begin{definition}
\textit{Направлением кривой} в точке \(z(t)\) будем называть \(\arg{z'(t)}\).
\end{definition}

\begin{definition}
\textit{Составную кривую \(\gamma\)} будем называть \textit{кусочно непрерывно дифференцируемой
(КНД)}, если \(\gamma\) состоит из конечного числа непрерывно-дифференцируемых кривых. Аналогично
вводится определение кусочно-гладкой функции.
\end{definition}

\begin{definition}
Открытое множество \(D \subset \oCo\) будем называть \textit{областью}, если множество \(D \setminus
\{\infty\}\) линейно связно, то есть любые две точки \(z_1, z_2 \in D \setminus \{\infty\}\) можно
соединить кусочнозвенной ломаной \(\gamma \subset D \setminus \{\infty\}\).
\end{definition}

\newpage
\begin{example}
\hfill \noindent
\par
\begin{minipage}[c]{0.45\textwidth}
    \centering
    \begin{tikzpicture}[>=stealth, scale=1.3]
        \fill[pattern=north west lines, pattern color=black] (0,0) circle (1);
        \draw[thick] (0,0) circle (1);

        \draw[->, thick] (-1.8, 0) -- (1.8, 0) node[right] {\(x\)};
        \draw[->, thick] (0, -1.8) -- (0, 1.8) node[above] {\(y\)};
        \node[below left] at (0,0) {\(0\)};

        \filldraw (-1,-1) circle (1.5pt) node[below left] {\(z_1\)};
        \filldraw (1,1) circle (1.5pt) node[above right] {\(z_2\)};

        \draw[thick] (-1,-1) -- (-1.4, 0.5) -- (0.5, 1.4) -- (1,1);
    \end{tikzpicture}
\end{minipage}
\hfill
\begin{minipage}[c]{0.5\textwidth}
Множество \(D = \{|z| > 1\} \cup \{\infty\}\) является линейно связным, так как любые две его точки
можно соединить ломаной. Также оно является открытым, что следует из определения. Значит, \(D\)~---
открыто. Отметим, что здесь и далее заштрихованная часть комплексной плоскости не является частью
множества.
\end{minipage}
\end{example}

\begin{definition}
Область \(D \subset \Co\) будем называть \textit{односвязной}, если любой простой контур из \(D\)
можно непрерывно деформировать в точку, оставаясь при этом в \(D\).
\end{definition}

\begin{example}
\par\hfill\break\noindent
\begin{minipage}[c]{0.45\linewidth}
    \centering
    \begin{tikzpicture}[scale=0.85]
        % Штриховка областей СНАРУЖИ полосы (y >= 1 и y <= -1)
        \pattern[pattern=north east lines, pattern color=gray!70]
            (-2.3, 1) rectangle (2.3, 2.0);
        \pattern[pattern=north east lines, pattern color=gray!70]
            (-2.3, -2.0) rectangle (2.3, -1);

        % Оси координат
        \draw[->, >=stealth, thick] (-2.5, 0) -- (2.5, 0) node[right] {\(x\)};
        \draw[->, >=stealth, thick] (0, -2.2) -- (0, 2.2) node[above] {\(y\)};

        % Пунктирные границы y = 1 и y = -1 (не принадлежат множеству)
        \draw[thick] (-2.3, 1) -- (2.3, 1) node[pos=0.9, above=-2pt] {\scriptsize \(i\)};
        \draw[thick] (-2.3, -1) -- (2.3, -1) node[pos=0.9, below=-2pt] {\scriptsize \(-i\)};

        % Метки на осях и начало координат
        \node[below left] at (0,0) {\(0\)};
        \filldraw (0, 1) circle (1.2pt);
        \filldraw (0, -1) circle (1.2pt);
    \end{tikzpicture}
\end{minipage}%
\begin{minipage}[c]{0.52\linewidth}
Множество \(D_1 = \{z \in \Co \colon |\im{z}| < 1\}\) является односвязным. Например, можно
деформировать к \(x = 0\), а после к нулю.
\end{minipage}

\vspace{2em}

\noindent
\begin{minipage}[c]{0.45\linewidth}
  \centering
  \begin{tikzpicture}[scale=0.85]
    % Штриховка ВНУТРИ круга радиуса 1
    \fill[pattern=north east lines, pattern color=gray!70]
      (0,0) circle (1);

    % Оси координат
    \draw[->, >=stealth, thick] (-2.5, 0) -- (2.5, 0) node[right] {\(x\)};
    \draw[->, >=stealth, thick] (0, -2.5) -- (0, 2.5) node[above] {\(y\)};

    % Пунктирная граница окружности |z| = 1
    \draw[thick] (0, 0) circle (1);

    % Метки на осях
    \node[below left] at (0,0) {\(0\)};
    \node[below right] at (1,0) {\scriptsize \(1\)};
    \node[above left] at (0,1) {\scriptsize \(i\)};

    \filldraw (1, 0) circle (1.2pt);
    \filldraw (0, 1) circle (1.2pt);
  \end{tikzpicture}
\end{minipage}%
\begin{minipage}[c]{0.52\linewidth}
Множество \(D_2 = \{z \in \Co \colon |z| > 1\}\) не является односвязным. Примером кривой, которую
невозможно непрерывно деформировать в точку, не выходя из \(D_2\) является окружность с центром в
нуле и радиусом, равным двум.
\end{minipage}
\end{example}


\Section{Последовательности и ряды}

\begin{definition}
Отображение \(\{z_n\} \colon \N \to \Co\) называется \textit{последовательностью комплексных чисел}.
\end{definition}

\begin{definition}
\textit{Пределом последовательности комплексных чисел \(\{z_n\}\)} является \(A = \lim\limits_{n \to
\infty} z_n\), если \(\forall \epsilon > 0 \ \exists N \in \N \colon \forall n \ge N \mapsto z_n \in
B_\epsilon(A)\). Если \(A \ne \infty\), то последовательность называется \textit{сходящейся}. В
случае сходящейся последовательности условие \(z_n \in B_\epsilon(A)\) эквивалентно \(|z-A| <
\epsilon\). Записи вида \(A = \lim\limits_{n \to \infty} z_n\) и \(z_n \xrightarrow{n \to \infty}
A\) эквивалентны.
\end{definition}

\begin{statement}\label{7:arithmetic_properties}
Если \(\lim\limits_{n \to \infty} z_n = A^*\), \(\lim\limits_{n \to \infty} w_n = B^*\), где \(A^*,
B^* \ne \infty\), то последовательности \(z_n + w_n, z_nw_n\) сходятся и их пределы равны \(A^* +
B^*, A^*B^*\) соответственно. Если кроме этого выполняется \(B^* \ne 0 \ \& \ w_n \ne 0 \forall n
\in \N\), то последовательность \(\dfrac{z_n}{w_n}\) тоже сходится и её предел равен
\(\dfrac{A^*}{B^*}\).
\end{statement}

\begin{proof}
Доказательство данного утверждения аналогично доказательствам из прошлых семестров.
\end{proof}

\begin{statement}\label{7:complex_to_real}
Пусть \(z_n = x_n + iy_n\)~--- последовательность комплексных чисел, где \(x_n, y_n\)~---
последовательности действительных чисел. Тогда условия \(\lim\limits_{n \to \infty} z_n = A + Bi\) и
\(x_n \xrightarrow{n \to \infty} A \ \& \ y_n \xrightarrow{n \to \infty} B\), где \(A, B\)~---
действительные числа, эквивалентны.
\end{statement}

\begin{proof}
Запишем ряд неравенств: \(\max \{|x_n -A|, |y_n - B|\} \le \sqrt{(x_n-A)^2 + (y_n - B)^2} = |z_n -
(A+Bi)| \le |x_n - A| + |y_n - B|\).

В случае первого условия \(\forall \epsilon \ \exists N_2 (\epsilon) = N_1(\epsilon) \colon \forall
n \ge N_2 \mapsto \max \{|x_n -A|, |y_n - B|\} \le |z_n - (A+Bi)| < \epsilon\). В случае второго
условия \(\forall \epsilon \ \exists N_1 = \max \{N_{2x}, N_{2y}\} (\epsilon) = N_1(\epsilon) \colon
\forall n \ge N_2 \mapsto|z_n - (A+Bi)| \le |x_n - A| + |y_n - B| < 2\epsilon\). Имеем определения
данных условий, а значит утверждение доказано.
\end{proof}

\begin{definition}
Сумму вида \(\sum\limits_{n=0}\limits^{\infty} z_n\) будем называть \textit{рядом}. Сумму \(S_n =
\sum\limits_{n=0}\limits^{k} z_n\) будем называть \textit{частной (частичной) суммой ряда}. Если
последовательность \(\{S_n\}\) сходится, то ряд называется \textit{сходящимся}. Если ряд
\(\sum\limits_{n=0}\limits^{\infty} |z_n|\) сходится, то ряд \(\sum\limits_{n=0}\limits^{\infty}
z_n\) называется \textit{абсолютно сходящимся}.
\end{definition}

\begin{note}
Из абсолютной сходимости ряда следует его сходимость. Обратное неверно.
\end{note}

\begin{statement}
Пусть \(z_n = x_n + iy_n\)~--- последовательность комплексных чисел, где \(x_n, y_n\)~---
последовательности действительных чисел. Тогда условия \(\sum\limits_{n=0}\limits^{\infty} z_n = S_1
+ S_2i\) и \(\sum\limits_{n=0}\limits^{\infty} x_n = S_1 \ \& \ \sum\limits_{n=0}\limits^{\infty}
y_n = S_2\), где \(S_1, S_2\)~--- действительные числа, эквивалентны.
\end{statement}

\begin{proof}
Доказательство следует из определения сходимости ряда и утверждения \ref{7:complex_to_real}.
\end{proof}

Вообще говоря, мы можем пользоваться некоторыми утверждениями из вещественного анализа, например
теоремой Больцано-Вейерштрасса или критерием Коши, так как они верны в пространстве \(\R^2\). При
этом в комплексном анализе отстутствуют операции сравнения, а значит нет аналогов утверждения о
монотонной ограниченной последовательности или <<теоремы о двух милиционерах>>.


\Section{Предел функции в точке. Непрерывность функции}

\begin{definition}
Пусть \(z_0, A \in \oCo\), причём \(z_0\)~--- предельная для множества \(E \subset \oCo\). \(A \in
\oCo\) называется \textit{пределом функции при \(z\) стремящемся к \(z_0\) по множеству \(E\)}, если
\(\forall \epsilon > 0 \ \exists \delta > 0 \colon \forall z \in (E \cap \mathring{B}_\delta (z_0))
\mapsto f(z) \in B_\epsilon(A)\) и обозначается \(\lim\limits_{z \to z_0, \ z \in E} f(z) = A\).
Если для некоторого \(\delta_0\) \(\mathring{B}_{\delta_0} (z_0) \subset E\), то в обозначении можно
опустить условие \(z \in E\).
\end{definition}

\begin{note}
Арифметические свойства предела функции в комплексном анализе подобны утверждению
\ref{7:arithmetic_properties}. Их доказательства также подобны доказательствам из прошлых семестров.
\end{note}

\begin{definition}
Пусть \(z_0 \in \Co\), причём \(z_0\)~--- предельная для множества \(E \subset \Co\) и принадлежит
ему. Если \(\lim\limits_{z \to z_0, \ z \in E} f(z) = f(z_0)\), то функция \(f\) называется
непрерывной в точке \(z_0\) по множеству \(E\). Если для некоторого \(\delta_0\)
\(\mathring{B}_{\delta_0} (z_0) \subset E\), то в обозначениях можно опустить условие \(z \in E\).
\end{definition}

\begin{example}
Функция
\begin{equation*}
    f(z) =
        \begin{cases}
            1/z, & z \neq 0, \infty; \\
            0, & z = \infty; \\
            \infty, & z = 0
        \end{cases}
\end{equation*}
является непрерывной на всей расширенной комплексной плоскости и является биекцией.
\end{example}

\begin{note}
Арифметические свойства непрерывных функции в комплексном анализе подобны утверждениям из прошлых
семестров. Их доказательства им подобны.
\end{note}

\begin{statement}
Если:
\begin{itemize}
\item \(f(z)\) определена в \(B_\delta (z_0)\) и непрерывна в \(z_0\);
\item \(f(z_0) = w_0 \in \Co\);
\item \(g(w)\) определена в \(B_\epsilon (w_0)\) и непрерывна в \(w_0\),
\end{itemize}
то \(g(f(z))\) определена в \(B_{\delta_1} (z_0)\) и непрерывна в \(z_0\).
\end{statement}

\begin{example}
\(\text{e}^{\cos^2{(\sin{z^3})}}\) непрерывна на всей комплексной плоскости.
\end{example}

\begin{statement}
Пусть \(z_0 = x_0 + iy_0\). Тогда \(f(z) = u(x, y) +iv(x,y)\) непрерывна в \(z_0\) тогда, и только
тогда, когда функции \(u(x, y)\) и \(v(x, y)\) непрерывны в \((x_0, y_0)\).
\end{statement}

\begin{proof}
Доказательство вытекает из ряда неравенств:
\begin{align}
    \max \{u(x,y) - u(x_0, y_0), \, v(x,y) - v(x_0, y_0)\}
    &\le |f(z) - f(z_0)| \le\\
    &\le |u(x,y) - u(x_0, y_0)| + |v(x,y) - v(x_0, y_0)|.
\end{align}
\end{proof}

\begin{example}
\(f(z) = \text{e}^z = \text{e}^x\cos{y} + i\text{e}^x\sin{y}\) непрерывна на всей комплексной
плоскости.
\end{example}


\Section{Дифференцирование по комплексной переменной}

\begin{definition}
Пусть \(f(z)\) определена в некоторой окрестности точки \(z_0\). Функцию \(f\) будем называть
дифференцируемой, если существует \(\lim\limits_{\Delta z \to 0}
\dfrac{f(z_0 + \Delta z) - f(z_0)}{\Delta z} = A \in \Co\). Число \(A = f'(z_0)\) будем наывать
производной функции \(f\) в точке \(z_0\). Обозначать будем \(f \in D(z_0)\).
\end{definition}

\begin{definition}\label{9:alt_dif}
Функцию \(f\) будем называть дифференцируемой, если существуют такое число \(A\) и функция
\(\alpha(\Delta z)\), определённая в некоторой окрестности нуля и непрерывная в нём, причём \(\alpha
(0) = 0\), и выполняется: \(f(z_0 + \Delta z) - f(z_0) = A\Delta z + \alpha (\Delta z) |\Delta z|\).
\end{definition}

\begin{note}
Из определения \ref{9:alt_dif} видно, что если функция дифференцируема, то она непрерывна. Обратное
неверно.
\end{note}

\begin{statement}
Пусть \(f(z) = u(x, y) + iv(x, y)\) определена в некоторой окрестности точки \(z_0\). Тогда
дифференцируемость функции \(f\) в точке \(z_0\) эквивалентна условию Коши-Римана:
\begin{equation*}
    \begin{cases}
        u(x, y), \ v(x,y) \in D(x_0, y_0), \\
        u_x (x_0, y_0) = v_y (x_0, y_0), \\
        u_y (x_0, y_0) = - v_x (x_0, y_0),
    \end{cases}
\end{equation*}
где \(f_t\)~---  обозначение частной производной функции \(f\) (\(u\) или \(v\)) по переменной \(t\)
(\(x\) или \(y\)).
\end{statement}
